\documentclass{article}

\usepackage{amsmath,amssymb}
\usepackage{xurl}
\usepackage[hidelinks]{hyperref}
\setlength{\emergencystretch}{2em}

% Shared commands for notes in docs/paper-gaps/.

\DeclareUnicodeCharacter{2080}{\ifmmode{}_{0}\else\textsubscript{0}\fi}
\DeclareUnicodeCharacter{2081}{\ifmmode{}_{1}\else\textsubscript{1}\fi}
\DeclareUnicodeCharacter{2082}{\ifmmode{}_{2}\else\textsubscript{2}\fi}
\DeclareUnicodeCharacter{2099}{\ifmmode{}_{n}\else\textsubscript{n}\fi}

\newcommand{\Lean}{\textsc{Lean}}
\ExplSyntaxOn
\NewDocumentCommand{\leanid}{m}{
  \ifmmode
    \textsf{\detokenize{#1}}
  \else
    \group_begin:
    \urlstyle{sf}
    \tl_set:Nn \l_tmpa_tl {#1}
    \tl_replace_all:Nnn \l_tmpa_tl {\_} {_}
    \exp_args:NV \path \l_tmpa_tl
    \group_end:
  \fi
}
\ExplSyntaxOff
\providecommand{\tr}{\operatorname{tr}}
\newcommand{\tnleanIssueBase}{https://github.com/LionSR/TNLean/issues/}
\newcommand{\tnleanPrBase}{https://github.com/LionSR/TNLean/pull/}
\newcommand{\tnleanDocBase}{https://sirui-lu.com/TNLean/docs/}
\newcommand{\ghissue}[1]{\href{\tnleanIssueBase#1}{GitHub issue \##1}}
\newcommand{\ghpr}[1]{\href{\tnleanPrBase#1}{PR \##1}}
\newcommand{\doclink}[1]{\href{\tnleanDocBase#1.html}{\path{#1}}}

% Dirac notation and the identity operator, as in blueprint/src/macros/common.tex.
% \providecommand keeps note-local definitions authoritative.
\usepackage{dsfont}
\providecommand{\ket}[1]{|#1\rangle}
\providecommand{\bra}[1]{\langle #1|}
\providecommand{\braket}[2]{\langle #1 | #2 \rangle}
\providecommand{\ketbra}[2]{|#1\rangle\!\langle #2|}
\providecommand{\Id}{\mathds{1}}
\providecommand{\one}{\mathds{1}}
\providecommand{\C}{\mathbb{C}}
\providecommand{\MN}[1]{M_{#1}(\C)}
\providecommand{\MD}{M_D(\C)}

% External referencing for paper sources.  The build helper writes sanitized
% auxiliary files under build/paper-aux/ and excludes duplicated source labels.
\usepackage{xr}

% Import a generated paper auxiliary file with a caller-supplied label prefix.
% The candidate paths support builds from docs/paper-gaps/, the repository root,
% and build/paper-aux/ itself.
\newcommand{\importpaperaux}[2]{%
  \IfFileExists{../../build/paper-aux/#2.aux}{%
    \externaldocument[#1]{../../build/paper-aux/#2}%
  }{%
    \IfFileExists{build/paper-aux/#2.aux}{%
      \externaldocument[#1]{build/paper-aux/#2}%
    }{%
      \IfFileExists{#2.aux}{%
        \externaldocument[#1]{#2}%
      }{}%
    }%
  }%
}

% General external reference macros
\newcommand{\paperref}[2]{\ref{#1-#2}}
\newcommand{\papereqref}[2]{(\ref{#1-#2})}

% CPSV16 source (1606.00608 / MPDO-22-12-17-2.tex) external document and shortcuts
\importpaperaux{cpsv16-}{1606.00608/MPDO-22-12-17-2-tnlean}
\newcommand{\cpsvref}[1]{\paperref{cpsv16}{#1}}
\newcommand{\cpsveqref}[1]{\papereqref{cpsv16}{#1}}

% Verdict marker: \gapnote{<kind>}{<status>}. Kind is the policy
% classification (clarification, local-correction, scope-restriction,
% unfaithful, false-source, open-gap); status is open, wip (elimination
% underway), resolved, or historical. Machine-read by the site index;
% prints a verdict line.
\newcommand{\gapnote}[2]{\par\noindent\textbf{Verdict:} #1 (#2).\par}


\title{A raw-weight counterexample to the pure-state ZCL equivalence in
Cirac--Perez-Garcia--Schuch--Verstraete 2017}
\author{}
\date{2026-07-14}

\begin{document}
\maketitle
\gapnote{false-source}{resolved}

\section{Source statement and failed inference}

The canonical form in Ref.~\cite{Cirac2016MPDO_arXiv} permits repeated copies
of a normal tensor with coefficients of modulus at most one, provided at least
one coefficient equals one. In the notation of the source equation
\path{eq:II_ABasicTensors}, the tensor has the form
\begin{align}
    A^i
        &= \bigoplus_{j=1}^g\bigoplus_{q=1}^{r_j}
        \mu_{j,q}X_{j,q}A_j^iX_{j,q}^{-1}.
    \label{eq:cpsv16_pure_zcl_raw_weight_counterexample_canonical_form}
\end{align}
Theorem~3.8 of Ref.~\cite{Cirac2016MPDO_arXiv}, labelled
\path{TheoremZCLPure} in the local source, asserts that a canonical-form tensor
has zero correlation length if and only if its transfer matrix is idempotent.

The converse argument at source lines~1248--1251 infers that
non-idempotence of the full transfer matrix $\mathbb E_A$ implies
non-idempotence of some normal-component transfer matrix
$\mathbb E_{j,j}$. This inference omits the copy coefficients. On the bond
block indexed by $(j,q;j',q')$, the full transfer matrix is
\begin{align}
    \left.\mathbb E_A\right|_{(j,q;j',q')}
        &= \mu_{j,q}\overline{\mu_{j',q'}}\mathbb E_{j,j'}.
    \label{eq:cpsv16_pure_zcl_raw_weight_counterexample_weighted_transfer_block}
\end{align}
The full transfer matrix can therefore have eigenvalues strictly between zero
and one even when every diagonal component $\mathbb E_{j,j}$ is idempotent.

\section{Counterexample}

Let the physical alphabet contain one letter, and let the scalar normal tensor
be
\begin{align}
    C^0 &= (1).
    \label{eq:cpsv16_pure_zcl_raw_weight_counterexample_scalar_normal_tensor}
\end{align}
Form the canonical tensor with two copies of $C^0$, carrying weights $1$ and
$1/2$. In sector-adapted coordinates,
\begin{align}
    A^0
        &= C^0\oplus\frac12 C^0
    \label{eq:cpsv16_pure_zcl_raw_weight_counterexample_two_copy_sum}\\
        &= \begin{pmatrix}
            1 & 0\\
            0 & \frac12
        \end{pmatrix}.
    \label{eq:cpsv16_pure_zcl_raw_weight_counterexample_diagonal_tensor}
\end{align}
This tensor has one BNT component and two raw copies. Its weights satisfy the
normalization in Ref.~\cite{Cirac2016MPDO_arXiv}: both moduli are at most one,
and one weight equals one.

For a chain of length $N$, the matrix-product-vector coefficient is
\begin{align}
    V^{(N)}(A)
        &= 1+2^{-N}.
    \label{eq:cpsv16_pure_zcl_raw_weight_counterexample_chain_coefficient}
\end{align}
This is also the coefficient obtained directly from the two-copy BNT
decomposition.

Local orthogonality is vacuous because there is only one BNT component. The
physical Hilbert space of every finite region is one-dimensional, so every
physical observable is scalar. Every two-region expectation is consequently
independent of separation. Thus $A$ satisfies the CID and local-orthogonality
conditions in Definition~3.6 of Ref.~\cite{Cirac2016MPDO_arXiv}.

The transfer map has the matrix units as eigenvectors. Its eigenvalues are
\begin{align}
    1,\qquad \frac12,\qquad \frac12,\qquad \frac14.
    \label{eq:cpsv16_pure_zcl_raw_weight_counterexample_transfer_spectrum}
\end{align}
Because $1/2$ and $1/4$ are not idempotent scalars, the transfer map is not
idempotent.

The same failure appears in the one-letter blocking equation required at a
renormalization fixed point. If a scalar $v$ satisfied the required relation,
then
\begin{align}
    (A^0)^2 &= vA^0.
    \label{eq:cpsv16_pure_zcl_raw_weight_counterexample_blocking_equation}
\end{align}
The first diagonal entry would give
\begin{align}
    1 &= v,
    \label{eq:cpsv16_pure_zcl_raw_weight_counterexample_first_entry}
\end{align}
whereas the second would then require
\begin{align}
    \frac14 &= \frac12,
    \label{eq:cpsv16_pure_zcl_raw_weight_counterexample_second_entry_contradiction}
\end{align}
which is false.

All four counterexample assertions have kernel-checked proofs.\footnote{Formal
conjunction:
\leanid{MPSTensor.halvedWeightTensor_counterexample_to_unrestricted_zcl_iff_rfp}
in \doclink{TNLean/MPS/RFP/BNTWeightCounterexample}.}

\section{Consequence for the theorem surface}

The unrestricted equivalence in Theorem~3.8 of
Ref.~\cite{Cirac2016MPDO_arXiv} is false under its stated raw-weight
canonical-form normalization. The tensor in
Eq.~\ref{eq:cpsv16_pure_zcl_raw_weight_counterexample_diagonal_tensor}
disproves the implication from zero correlation length to transfer
idempotence. It does not settle the converse implication for adjacent physical
regions; the existing formal result in that direction assumes that both
complementary gaps are positive.

A corrected theorem must remove the copy-weight eigenvalues from its
conclusion. One option is to pass to a canonical representative containing one
unit-weight copy of each BNT component. Another is to state an asymptotic
property of the limiting transfer projector rather than literal idempotence of
the raw tensor's transfer map. Merely requiring unit-modulus copy weights does
not suffice because distinct phases can also destroy idempotence. Any
replacement must specify how it treats repeated copies and their phases.

The first option is now proved. For the direct sum containing one unit-weight
copy of each distinct BNT component, positive-gap physical CID together with
BNT local orthogonality is equivalent to transfer idempotence. The reverse
implication uses arbitrary sector-supported rank-one observables and
nondegeneracy of the trace pairing. It therefore avoids the false inference at
source line~1250 that non-idempotence supplies a nonzero subleading eigenvalue.
This representative-level theorem does not apply to
Eq.~\ref{eq:cpsv16_pure_zcl_raw_weight_counterexample_diagonal_tensor} and
does not alter the counterexample verdict.

The target tracked by \ghissue{2597} should not retain the unrestricted
biconditional. Its replacement must distinguish the false raw-tensor statement
from any theorem about a selected representative of the generated
matrix-product-vector family.

The verdict is a source counterexample: under the canonical-form normalization
stated in Ref.~\cite{Cirac2016MPDO_arXiv}, zero correlation length does not
imply transfer idempotence. This is not an open formalization gap.

\bibliographystyle{alpha}
\bibliography{references}

\end{document}
